\section{Methodology}
\label{methodology}

\subsection{Framework Overview}
\label{sec:method-overview}

The proposed method takes natural-language requirements as input and produces an architecture-aware alignment artifact.
The artifact contains architectural concerns, candidate components, responsibility assignments, architecture constraints, and trace links from requirements to architectural decisions.

\subsection{Multi-Agent Architecture Alignment Workflow}
\label{sec:method-workflow}

The workflow uses three coordinated roles rather than a large agent hierarchy.
Each role operates on the same requirement set and exchanges structured intermediate artifacts for review.

\subsubsection{Analyst Agent}
\label{sec:method-analyst-agent}

The analyst agent extracts functional requirements, quality attributes, constraints, domain terms, and architecturally significant concerns from natural-language requirements.

\subsubsection{Architecture Agent}
\label{sec:method-architecture-agent}

The architecture agent proposes candidate components, connectors, responsibilities, and design constraints based on the extracted architectural concerns and the available architecture knowledge.

\subsubsection{Reviewer Agent}
\label{sec:method-reviewer-agent}

The reviewer agent checks whether the proposed mapping is complete, consistent, and traceable.
It identifies unmapped requirements, weak responsibility assignments, and architecture decisions that lack requirement evidence.

\subsection{Requirement Semantic Analysis}
\label{sec:method-semantic-analysis}

This step structures requirements into architecture-relevant categories.

\subsubsection{Functional Requirements}
\label{sec:method-functional-requirements}

Functional requirements are analyzed to identify system capabilities, user-facing behaviors, and candidate responsibilities.

\subsubsection{Constraints}
\label{sec:method-constraints}

Constraints are extracted as limitations on technology, integration, compliance, deployment, data handling, or operational behavior.

\subsubsection{Quality Attributes}
\label{sec:method-quality-attributes}

Quality attributes such as performance, availability, scalability, maintainability, security, and usability are represented as architecture-driving concerns.

\subsubsection{Architectural Concerns}
\label{sec:method-architectural-concerns}

Architectural concerns summarize the design implications that should influence component boundaries, communication patterns, and architecture constraints.

\subsection{Knowledge-Guided Architecture Reasoning}
\label{sec:method-knowledge}

Architecture reasoning is guided by reusable knowledge rather than relying only on free-form generation.

\subsubsection{Architecture Knowledge}
\label{sec:method-architecture-knowledge}

The knowledge source contains architecture concepts, common component types, connector semantics, quality-attribute tactics, and design-decision criteria.

\subsubsection{Design Patterns}
\label{sec:method-design-patterns}

Design patterns are used as reasoning aids when requirements imply recurring responsibility structures or interaction styles.

\subsubsection{Architecture Templates}
\label{sec:method-architecture-templates}

Architecture templates provide lightweight reference structures for common system families and help standardize mapping outputs.

\subsubsection{Prompt Strategies}
\label{sec:method-prompt-strategies}

Prompt strategies constrain the agents to produce structured mappings, cite requirement evidence, and separate assumptions from supported architecture decisions.

\subsection{Requirement-to-Architecture Mapping}
\label{sec:method-mapping}

This step is the central output of the method.
It links requirement-level evidence to architectural elements and design constraints.

\subsubsection{Component Identification}
\label{sec:method-component-identification}

Candidate components are identified from functional responsibilities, domain concepts, external actors, data ownership, and quality-attribute drivers.

\subsubsection{Responsibility Allocation}
\label{sec:method-responsibility-allocation}

Responsibilities are assigned to components with explicit references to the requirement statements that justify each assignment.

\subsubsection{Architecture Constraint Alignment}
\label{sec:method-constraint-alignment}

Architecture constraints are checked against extracted requirement constraints and quality attributes to reduce requirement-design inconsistency.

\subsubsection{Traceability Construction}
\label{sec:method-traceability}

Trace links connect requirements, architectural concerns, components, responsibilities, constraints, and reviewer comments.

\subsection{Human-in-the-Loop Interaction}
\label{sec:method-hitl}

Human reviewers validate ambiguous mappings, approve assumptions with architectural impact, and revise agent outputs when requirement evidence is insufficient.
The interaction is limited to requirement-to-architecture alignment decisions.
